\documentclass[../../main.tex]{subfiles}
\graphicspath{{../../image/}{../../figure/}}
\begin{document}

\section{不可压缩性}\label{section:ch02-不可压缩性}

许多流体，甚至某些固体（例如橡胶），可较容易地被拉伸或剪切，但在受到会改变其体积的变形时却表现出不成比例的高刚度。连续介质物理将这类材料视为不能承受任何体积变化，于是密度 $\rho$ 沿质点轨迹保持常数。不可压缩性条件
\begin{align}
\det F=1,\qquad\mathrm{tr}\,D=\mathrm{div}\,v^{\top}=0,
\label{align:ch02-2-7-1}
\end{align}
以 Lagrange 或 Euler 坐标写出，作为一个运动学约束附加到平衡律方程组。作为回报，应力张量分解为两部分：
\begin{align}
S=-p(F^{-1})^{\top}+\hat{S},\qquad T=-pI+\hat{T},
\label{align:ch02-2-7-2}
\end{align}
其中 $\hat{S}$ 或 $\hat{T}$ 称为\emph{附加应力}，如前一样由热力学过程通过本构关系确定；而另一项表示一个\emph{静水压强}，它不由本构关系指定，而须与热力学过程一起，通过在运动学约束\eqref{align:ch02-2-7-1} 之下求解质量、动量和能量平衡律方程组来确定。

静水压强的显著性质是它在等容变形下不做功。为用热力学第二定律来激励\eqref{align:ch02-2-7-2}，考察一个不可压缩热弹性材料，其 $\varepsilon,\theta,Q$ 的本构方程如\eqref{align:ch02-2-5-1}，但仅对满足 $\det F=1$ 的 $F$ 有定义，而 $S$ 不加指定。耗散不等式再次给出\eqref{align:ch02-2-5-2}，其中 $S$ 被替换、$\partial_{F}\hat{\varepsilon}$ 被替换为流形 $\det F=1$ 上的切向导数 $\partial^{\tau}_{F}\hat{\varepsilon}$，并且 $\dot{F}$ 被限制在子空间
\begin{align}
\mathrm{tr}\,[(F^{-1})^{\top}\dot{F}^{\top}]=\mathrm{tr}\,[(F^{*})^{\top}\dot{F}^{\top}]=\mathrm{tr}\,[(\partial_{F}\det F)\dot{F}^{\top}]=\dot{\det F}=0
\label{align:ch02-2-7-3}
\end{align}
内。因此，$\mathrm{tr}\,[(\rho_{0}\partial^{\tau}_{F}\hat{\varepsilon}-S)\dot{F}^{\top}]\leqslant 0$ 对所有满足\eqref{align:ch02-2-7-3} 的 $\dot{F}$ 成立，当且仅当
\begin{align}
S=-p(F^{-1})^{\top}+\rho_{0}\partial^{\tau}_{F}\hat{\varepsilon}(F,s),
\label{align:ch02-2-7-4}
\end{align}
对某个标量 $p$ 成立。

在不可压缩 Newton 流体中，应力仍由\eqref{align:ch02-2-6-16} 给出，但其中 $\rho$ 为常数，而 $p(\rho,s)$ 被替换为未定的静水压强 $p$。当不可压缩流体是无粘时，整个应力张量都被未定静水压强所取代。

\end{document}
